\section{Fulton--MacPherson Calculus: Orientation and Residue Formulas}
\label{app:FM_Stokes}

The Fulton--MacPherson calculus used in Section~\ref{sec:FM_calculus}
rests on four explicit ingredients: the boundary orientation convention,
the logarithmic residue formula, Arnold cancellation at codimension-two
corners, and the form-degree proof of cubic Landau--Ginzburg truncation.

\subsection{Orientation Conventions and Sign Rules}

\begin{proposition}[Orientation formula]
\label{prop:orientation_formula}
On the FM compactification $\FM_k(\C)$ with boundary divisor $D_S$, the orientation is determined by
\[
\text{or}(\FM_k(\C)) = (-d\varepsilon_S) \wedge \text{or}(D_S),
\]
where $\varepsilon_S \geq 0$ is the inward-pointing radial parameter
($\varepsilon_S = 0$ at the boundary, $\varepsilon_S > 0$ in the interior)
and the outward normal is $-\partial_{\varepsilon_S}$
(Convention~\ref{conv:outward-normal}).

For logarithmic forms $\omega \in \Omega^{k-1}_{\log}(\FM_k(\C))$, the residue operator $\Res_{D_S}$ extracts the coefficient of $d\log\varepsilon_S$:
\[
\Res_{D_S}(\omega) = \beta\big|_{\varepsilon_S = 0}, \quad \text{where } \omega = d\log\varepsilon_S \wedge \beta + \alpha,
\]
with $\alpha, \beta$ smooth at $\varepsilon_S = 0$. Equivalently, $\Res_{D_S}(\omega) = \lim_{\varepsilon_S \to 0}\bigl(\varepsilon_S \cdot \iota_{\partial_{\varepsilon_S}} \omega\bigr)\big|_{D_S}$. Note that the naive contraction $\iota_{\partial_{\varepsilon_S}} \omega\big|_{D_S}$ (without the factor $\varepsilon_S$) diverges for logarithmic forms and must not be used.
The sign in Stokes' theorem is fixed by the convention that boundary orientations are induced from the manifold orientation via the outward normal first rule.
\end{proposition}

\subsection{Explicit Computation for $k=3$: All Boundary Divisors}
\label{subsec:k3_complete}

The $n=3$ $A_\infty$ identity follows by computing the residue at each codimension-1 boundary divisor of $\FM_3(\C)$ and assembling the contributions via Stokes' theorem.

\subsubsection{Setup and Coordinates}

After quotienting \(\C^3\) by translations only,
\(\FM_3(\C)=\FM_3^{\mathrm{tr}}(\C)\) has real dimension \(2(3-1)=4\).
Use quotient coordinates \(w_1=z_1-z_3\), \(w_2=z_2-z_3\), so the
complex orientation form is
\[
\omega_{\FM} = \tfrac{i}{2}\, dw_1 \wedge d\bar{w}_1 \wedge \tfrac{i}{2}\, dw_2 \wedge d\bar{w}_2.
\]
The weight form for three inputs $a, b, c$ is the closed logarithmic form
\[
\omega_3(a,b,c) = \Omega(w_1, w_2) \cdot a(z_1)\, b(z_2)\, c(z_3),
\]
where $\Omega$ is built from propagators and has logarithmic singularities along each boundary divisor. The four codimension-1 boundary divisors are $D_{\{1,2\}}$, $D_{\{1,3\}}$, $D_{\{2,3\}}$, and $D_{\{1,2,3\}}$.

\subsubsection{Divisor $D_{\{1,2\}}$: Points 1 and 2 Collide}

\begin{computation}[$D_{\{1,2\}}$ residue]
\label{comp:D12}
Near $D_{\{1,2\}}$, introduce FM coordinates:
\begin{itemize}
\item $\varepsilon_{12} \geq 0$: radial separation $|z_1 - z_2|$;
\item $u_{12} = \tfrac{1}{2}(z_1 + z_2) \in \C$: cluster center;
\item $\theta_{12} \in S^1$: angle $\arg(z_1 - z_2)$;
\item $z_3 \in \C$: position of the third point.
\end{itemize}
In these coordinates, $z_1 - z_2 = \varepsilon_{12}\, e^{i\theta_{12}}$ and the quotient coordinates become
\[
w_1 = u_{12} - z_3 + \tfrac{1}{2}\varepsilon_{12}\, e^{i\theta_{12}}, \qquad
w_2 = u_{12} - z_3 - \tfrac{1}{2}\varepsilon_{12}\, e^{i\theta_{12}}.
\]
The propagator factor $1/(z_1 - z_2) = 1/(\varepsilon_{12}\, e^{i\theta_{12}})$ produces a logarithmic singularity:
\[
\frac{d(z_1 - z_2)}{z_1 - z_2} = d\log\varepsilon_{12} + i\, d\theta_{12}.
\]
The weight form thus decomposes as
\[
\omega_3 = d\log\varepsilon_{12} \wedge \beta_{12} + \alpha_{12},
\]
where $\beta_{12}$ and $\alpha_{12}$ are smooth at $\varepsilon_{12} = 0$. The residue is
\[
\Res_{D_{\{1,2\}}}(\omega_3) = \beta_{12}\big|_{\varepsilon_{12}=0}.
\]
On
\[
D_{\{1,2\}}\cong
\FM_2^{\mathrm{tr}}(\C)\times\FM_2^{\mathrm{red}}(\C),
\]
the first factor records the pair \((u_{12},z_3)\) and the second
records the relative angle \(\theta_{12}\) within the cluster. By the
factorization property (Theorem~\ref{thm:cluster_factorization}), this
residue equals the wedge product
\(\omega_2^{\text{inner}}\wedge\omega_2^{\text{outer}}\), where:
\begin{itemize}
\item $\omega_2^{\text{inner}}$ is the weight form for $m_2(a, b)$ with spectral parameter $\lambda = z_1 - z_2$;
\item $\omega_2^{\text{outer}}$ is the weight form for $m_2(\,\cdot\,, c)$ with the cluster output at position $u_{12}$ and parameter $\mu = u_{12} - z_3$.
\end{itemize}
The spectral substitution principle dictates that the outer parameter is evaluated at $\mu = u_{12} - z_3$, which after setting $\varepsilon_{12} = 0$ gives $\mu = w_1 = w_2 = u_{12} - z_3$. The composed operation is
\[
\Res_{D_{\{1,2\}}}(\omega_3) \;\longleftrightarrow\; m_2\bigl(m_2(a, b;\, \lambda),\, c;\, \mu\bigr).
\]
The incidence sign is \(\epsilon_{D_{\{1,2\}}}=+1\),
computed from the orientation convention
\(\operatorname{or}(\FM_3(\C))=(-d\varepsilon_{12})
\wedge\operatorname{or}(D_{\{1,2\}})\) and the comparison with
the product orientation on
\(\FM_2(\C)\times\FM_2^{\mathrm{red}}(\C)\).
\end{computation}

\subsubsection{Divisor $D_{\{2,3\}}$: Points 2 and 3 Collide}

\begin{computation}[$D_{\{2,3\}}$ residue]
\label{comp:D23}
Near $D_{\{2,3\}}$, introduce FM coordinates:
\begin{itemize}
\item $\varepsilon_{23} \geq 0$: radial separation $|z_2 - z_3|$;
\item $u_{23} = \tfrac{1}{2}(z_2 + z_3) \in \C$: cluster center;
\item $\theta_{23} \in S^1$: angle $\arg(z_2 - z_3)$;
\item $z_1 \in \C$: position of the first point.
\end{itemize}
In these coordinates, $z_2 - z_3 = \varepsilon_{23}\, e^{i\theta_{23}}$, and the quotient coordinates are
\[
w_1 = z_1 - u_{23} + \tfrac{1}{2}\varepsilon_{23}\, e^{i\theta_{23}}, \qquad
w_2 = \varepsilon_{23}\, e^{i\theta_{23}}.
\]
The propagator $1/(z_2 - z_3)$ again produces $d\log\varepsilon_{23} + i\, d\theta_{23}$, and the residue at $\varepsilon_{23} = 0$ gives
\[
\Res_{D_{\{2,3\}}}(\omega_3) = \beta_{23}\big|_{\varepsilon_{23}=0}.
\]
On
\[
D_{\{2,3\}}\cong
\FM_2^{\mathrm{tr}}(\C)\times\FM_2^{\mathrm{red}}(\C),
\]
the factorization gives:
\begin{itemize}
\item Inner: $m_2(b, c;\, \mu)$ with $\mu = z_2 - z_3$ (the spectral parameter for the colliding pair);
\item Outer: $m_2(a,\, \cdot\,;\, \lambda)$ with $\lambda = z_1 - u_{23}$.
\end{itemize}
After setting $\varepsilon_{23} = 0$, the outer spectral parameter becomes $\lambda = z_1 - z_3 = w_1$. This is the same as $\lambda_1 + \lambda_2$ in the original notation (where $\lambda_1 = z_1 - z_2$ and $\lambda_2 = z_2 - z_3 = \mu$), consistent with the spectral substitution rule: the inner block $\{2,3\}$ is replaced by the sum $\Lambda_{\{2,3\}} = \lambda_2 \to \mu$, and the outer parameter for slot 1 relative to the collapsed slot is $\lambda_1 + \mu$. The composed operation is
\[
\Res_{D_{\{2,3\}}}(\omega_3) \;\longleftrightarrow\; (-1)^{|a|}\, m_2\bigl(a,\, m_2(b,c;\,\mu);\, \lambda + \mu\bigr).
\]
The Koszul sign $(-1)^{|a|}$ arises from passing the degree-$(-1)$ operation $m_2$ past $a$ in the outer composition (the standard sign $(-1)^{\epsilon(s,j)}$ with $s=1$, $j=2$: $\epsilon(1,2) = (2-1)\cdot |a| = |a|$).

The incidence sign is \(\epsilon_{D_{\{2,3\}}}=+1\): in the
change of coordinates from \((w_1,w_2)\) to
\((z_1,\varepsilon_{23},\theta_{23},u_{23})\), the transposition
parity and the outward-normal factor give the same product
orientation on
\(\FM_2(\C)\times\FM_2^{\mathrm{red}}(\C)\).
\end{computation}

\subsubsection{Divisor $D_{\{1,3\}}$: Points 1 and 3 Collide}

\begin{computation}[$D_{\{1,3\}}$ residue]
\label{comp:D13}
Near $D_{\{1,3\}}$, introduce FM coordinates:
\begin{itemize}
\item $\varepsilon_{13} \geq 0$: radial separation $|z_1 - z_3|$;
\item $u_{13} = \tfrac{1}{2}(z_1 + z_3) \in \C$: cluster center;
\item $\theta_{13} \in S^1$: angle $\arg(z_1 - z_3)$;
\item $z_2 \in \C$: position of the second point.
\end{itemize}
In these coordinates, $z_1 - z_3 = \varepsilon_{13}\, e^{i\theta_{13}}$ and $w_1 = \varepsilon_{13}\, e^{i\theta_{13}}$. The propagator $1/(z_1 - z_3)$ produces a logarithmic singularity, and the residue at $\varepsilon_{13} = 0$ factorizes
\[
D_{\{1,3\}}\cong
\FM_2^{\mathrm{tr}}(\C)\times\FM_2^{\mathrm{red}}(\C)
\]
into:
\begin{itemize}
\item Inner: $m_2(a, c)$ with the pair $(z_1, z_3)$ colliding;
\item Outer: the result composed with $b$ at position $z_2$.
\end{itemize}
However, \(\{1,3\}\) is a \emph{non-consecutive} subset: in the
ordered sequence of inputs \((a_1,a_2,a_3)=(a,b,c)\), the subset
\(\{1,3\}\) skips input \(2\).  The divisor \(D_{\{1,3\}}\) is present
in the closed chiral FM space.  It contributes zero only in the open
planar \(E_1\)-ordered projection used for the tree-level
\(A_\infty\) identity.  In that projection the HT propagator
\(K(t,z)=\Theta(t)/(2\pi z)\) and the chamber \(t_1>t_2>t_3\) force
planar edges to be consecutive.  The contraction \(z_1\to z_3\) with
\(z_2\) outside the cluster has no internal planar edge producing
\(d\log\varepsilon_{13}\).  The restriction of \(\omega_3\) to
\(D_{\{1,3\}}\) is regular in \(\varepsilon_{13}\).  Therefore, for
the planar tree-level weight form,
\[
\Res_{D_{\{1,3\}}}(\omega_3) = 0.
\]
The incidence sign is \(\epsilon_{D_{\{1,3\}}}=-1\): the radial
coordinate \(d\varepsilon_{13}\) is already in the first normal
slot, so the sign is the outward-normal factor.  The residue
vanishes, hence this face contributes zero.
\end{computation}

\subsubsection{Divisor $D_{\{1,2,3\}}$: All Three Points Collide}

\begin{computation}[$D_{\{1,2,3\}}$ contribution]
\label{comp:D123}
The divisor \(D_{\{1,2,3\}}\) corresponds to all three points
approaching a common limit simultaneously. Then
\[
D_{\{1,2,3\}} \cong
\FM_1^{\mathrm{tr}}(\C)\times\FM_3^{\mathrm{red}}(\C)
\cong \{\mathrm{pt}\}\times\FM_3^{\mathrm{red}}(\C)
\cong \FM_3^{\mathrm{red}}(\C),
\]
since \(\FM_1^{\mathrm{tr}}(\C)\) is a single point. Near
\(D_{\{1,2,3\}}\), the FM coordinates are:
\begin{itemize}
\item $\varepsilon_{123} \geq 0$: overall scale of the three-point cluster;
\item Internal coordinates on $\FM_3(\C)$: the relative configuration of all three points.
\end{itemize}
The factorization at this boundary gives two types of composed operations:
\begin{itemize}
\item $m_1 \circ m_3$: the differential $Q = m_1$ acts on the output of $m_3(a,b,c)$, giving $m_1(m_3(a,b,c))$;
\item $m_3 \circ (m_1 \otimes \id \otimes \id)$ and its permutations: $m_3$ with $m_1 = Q$ applied to one input.
\end{itemize}
These contribute the terms
\begin{align*}
&m_1(m_3(a,b,c)) + m_3(m_1(a), b, c) + (-1)^{|a|}\, m_3(a, m_1(b), c) + (-1)^{|a|+|b|}\, m_3(a, b, m_1(c)).
\end{align*}
The signs are the standard Koszul signs $(-1)^{\epsilon(s,1)} = (-1)^{0 \cdot |a_1 \cdots a_s|}$ for $j=1$: passing the degree-$0$ operation $m_1$ past the preceding inputs gives $(-1)^{(1-1)(|a_1|+\cdots+|a_s|)} = (-1)^0 = 1$ when $m_1$ hits $a_1$, and $(-1)^{|a|}$ when it hits $a_2$ (since then $m_1$ is being inserted at position $s=1$ in $m_3$, and the relevant sign in the $A_\infty$ identity~\eqref{eq:ainfty-relation-raw} is $(-1)^{\epsilon(s,j)}$ with $j=1$, $s$ the number of preceding inputs). More precisely, the Stasheff sign convention gives $\epsilon(s,1) = (1-1)(|a_1|+\cdots+|a_s|) = 0$ for all $s$, but the suspended degree conventions for $m_3$ of degree $1-3 = -2$ introduce the signs $(-1)^{|a|}$ and $(-1)^{|a|+|b|}$ via the Leibniz rule for the differential acting on the inputs of $m_3$.
\end{computation}

\subsubsection{Assembly: The $n=3$ $A_\infty$ Identity}

\begin{proposition}[$k=3$ Stasheff identity]
\label{prop:k3_identity}
Summing over all boundary divisors via Stokes' theorem (Theorem~\ref{thm:Stokes_FM}), and using $d\omega_3 = 0$ in the interior $\Conf_3(\C)$:
\[
0 = \int_{\Gamma_3} d\omega_3 =
\sum_{\substack{S \subseteq \{1,2,3\} \\ |S| \geq 2}}
\epsilon_{D_S}\int_{\Gamma_3 \cap D_S}\Res_{D_S}(\omega_3).
\]
The four contributions are:
\begin{enumerate}[label=(\roman*)]
\item \(D_{\{1,2\}}\), \(\epsilon_{D_{\{1,2\}}}=+1\):
contributes \(m_2(m_2(a,b;\lambda),c;\mu)\);
\item \(D_{\{2,3\}}\), \(\epsilon_{D_{\{2,3\}}}=+1\):
contributes \((-1)^{|a|}m_2(a,m_2(b,c;\mu);\lambda+\mu)\);
\item \(D_{\{1,3\}}\), \(\epsilon_{D_{\{1,3\}}}=-1\):
contributes \(0\) because the planar residue vanishes;
\item $D_{\{1,2,3\}}$: contributes $m_1(m_3(a,b,c)) + m_3(m_1(a),b,c) + (-1)^{|a|}\, m_3(a,m_1(b),c) + (-1)^{|a|+|b|}\, m_3(a,b,m_1(c))$.
\end{enumerate}
Setting the total to zero reproduces exactly the $n=3$ $A_\infty$ identity~\eqref{eq:ainfty-relation-raw}:
\begin{align*}
0 &= m_1(m_3(a,b,c)) + m_3(m_1(a),b,c) + (-1)^{|a|}\, m_3(a,m_1(b),c) + (-1)^{|a|+|b|}\, m_3(a,b,m_1(c)) \\
&\quad + m_2(m_2(a,b),c) + (-1)^{|a|}\, m_2(a,m_2(b,c)),
\end{align*}
where in the last line we have suppressed the spectral parameters, with the understanding that each $m_2$ composition carries the appropriate spectral substitution as computed above. \qed
\end{proposition}

\subsection{Explicit Computation for $k=4$}
\label{subsec:k4_computation}

The full boundary analysis for $\FM_4(\C)$ verifies the $n=4$ $A_\infty$ identity.

\subsubsection{Setup}

The translation-reduced FM compactification
\(\FM_4(\C)=\FM_4^{\mathrm{tr}}(\C)\) has real dimension \(2(4-1)=6\).
Use quotient coordinates \(w_i=z_i-z_4\) for \(i=1,2,3\). The boundary
divisors \(D_S\) with \(|S|\ge2\) are:
\begin{itemize}
\item \textbf{Two-element subsets} (\(|S|=2\)):
\(D_{\{1,2\}}, D_{\{1,3\}}, D_{\{1,4\}}, D_{\{2,3\}},
D_{\{2,4\}}, D_{\{3,4\}}\). Each gives
\(D_S\cong\FM_3^{\mathrm{tr}}(\C)\times
\FM_2^{\mathrm{red}}(\C)\).
\item \textbf{Three-element subsets} (\(|S|=3\)):
\(D_{\{1,2,3\}}, D_{\{1,2,4\}}, D_{\{1,3,4\}},
D_{\{2,3,4\}}\). Each gives
\(D_S\cong\FM_2^{\mathrm{tr}}(\C)\times
\FM_3^{\mathrm{red}}(\C)\).
\item \textbf{Four-element subset} (\(|S|=4\)):
\(D_{\{1,2,3,4\}}\cong
\FM_1^{\mathrm{tr}}(\C)\times\FM_4^{\mathrm{red}}(\C)
\cong\FM_4^{\mathrm{red}}(\C)\).
\end{itemize}

\subsubsection{Non-Vanishing Consecutive Divisors}

After projection to the open planar \(E_1\)-ordered chain
(Remark~\ref{rem:nonconsecutive_vanishing}), only consecutive subsets
contribute non-zero residues to the tree-level \(A_\infty\) identity.
For ordered inputs \((a_1,a_2,a_3,a_4)\), the consecutive subsets of
size \(\geq2\) are:
\[
\{1,2\},\quad \{2,3\},\quad \{3,4\},\quad \{1,2,3\},\quad \{2,3,4\},\quad \{1,2,3,4\}.
\]
Non-consecutive subsets
(\(\{1,3\}\), \(\{1,4\}\), \(\{2,4\}\),
\(\{1,3,4\}\), \(\{1,2,4\}\)) are still boundary divisors in the
closed chiral FM space.  Their residues vanish only after the planar
open projection, by the same argument as Computation~\ref{comp:D13}.

\subsubsection{Two-Element Consecutive Divisors}

\begin{computation}[$D_{\{1,2\}}$ in $\FM_4(\C)$]
Points 1 and 2 collide while 3 and 4 remain separate. The factorization is
\[
D_{\{1,2\}} \cong \FM_3(\C) \times \FM_2^{\mathrm{red}}(\C),
\]
where $\FM_3(\C)$ records positions of the cluster center $u_{12}$ plus points $z_3, z_4$ (three points total), and $\FM_2^{\mathrm{red}}(\C)$ records the relative position within the cluster. The residue gives
\[
\Res_{D_{\{1,2\}}}(\omega_4) \;\longleftrightarrow\; m_3\bigl(m_2(a_1, a_2;\, \lambda_1),\, a_3,\, a_4;\, \Lambda_{\{1,2\}}, \lambda_2, \lambda_3\bigr),
\]
with effective outer parameter $\Lambda_{\{1,2\}} = \lambda_1$ for the collapsed block, and the remaining parameters $\lambda_2, \lambda_3$ unchanged. The Koszul sign is $(-1)^{\epsilon(0,2)} = (-1)^{(2-1)\cdot 0} = +1$.
\end{computation}

\begin{computation}[$D_{\{2,3\}}$ in $\FM_4(\C)$]
Points 2 and 3 collide. The factorization
$D_{\{2,3\}} \cong \FM_3(\C) \times \FM_2^{\mathrm{red}}(\C)$ gives
\[
\Res_{D_{\{2,3\}}}(\omega_4) \;\longleftrightarrow\; (-1)^{|a_1|}\, m_3\bigl(a_1,\, m_2(a_2, a_3;\, \lambda_2),\, a_4;\, \lambda_1 + \Lambda_{\{2,3\}}, \lambda_3\bigr),
\]
where $\Lambda_{\{2,3\}} = \lambda_2$ is the effective parameter for the collapsed block, and the outer parameter for slot 1 becomes $\lambda_1 + \lambda_2$ (the distance from $z_1$ to the cluster center, by spectral substitution). The Koszul sign is $(-1)^{\epsilon(1,2)} = (-1)^{|a_1|}$.
\end{computation}

\begin{computation}[$D_{\{3,4\}}$ in $\FM_4(\C)$]
Points 3 and 4 collide. The factorization gives
\[
\Res_{D_{\{3,4\}}}(\omega_4) \;\longleftrightarrow\; (-1)^{|a_1|+|a_2|}\, m_3\bigl(a_1,\, a_2,\, m_2(a_3, a_4;\, \lambda_3);\, \lambda_1, \lambda_2 + \Lambda_{\{3,4\}}\bigr),
\]
where $\Lambda_{\{3,4\}} = \lambda_3$ and the Koszul sign is $(-1)^{\epsilon(2,2)} = (-1)^{|a_1|+|a_2|}$.
\end{computation}

\subsubsection{Three-Element Consecutive Divisors}

\begin{computation}[$D_{\{1,2,3\}}$ in $\FM_4(\C)$]
Points 1, 2, 3 collide while point 4 remains separate. The factorization is
\[
D_{\{1,2,3\}} \cong \FM_2(\C) \times \FM_3^{\mathrm{red}}(\C),
\]
where $\FM_2(\C)$ records the pair (cluster center $u_{123}$, point $z_4$) and $\FM_3^{\mathrm{red}}(\C)$ records relative positions within the cluster. The residue gives
\[
\Res_{D_{\{1,2,3\}}}(\omega_4) \;\longleftrightarrow\; m_2\bigl(m_3(a_1, a_2, a_3;\, \lambda_1, \lambda_2),\, a_4;\, \Lambda_{\{1,2,3\}}, \lambda_3\bigr),
\]
with $\Lambda_{\{1,2,3\}} = \lambda_1 + \lambda_2$ and Koszul sign $(-1)^{\epsilon(0,3)} = (-1)^{(3-1)\cdot 0} = +1$.
\end{computation}

\begin{computation}[$D_{\{2,3,4\}}$ in $\FM_4(\C)$]
Points 2, 3, 4 collide while point 1 remains separate. The factorization is
\[
D_{\{2,3,4\}} \cong \FM_2(\C) \times \FM_3^{\mathrm{red}}(\C),
\]
giving
\[
\Res_{D_{\{2,3,4\}}}(\omega_4) \;\longleftrightarrow\; (-1)^{|a_1|}\, m_2\bigl(a_1,\, m_3(a_2, a_3, a_4;\, \lambda_2, \lambda_3);\, \lambda_1 + \Lambda_{\{2,3,4\}}\bigr),
\]
with $\Lambda_{\{2,3,4\}}=\lambda_2+\lambda_3$.  In the
desuspended coordinate convention of
\eqref{eq:ainfty-relation-raw}, the insertion sign is
$(-1)^{\epsilon(1,3)}=(-1)^{|a_1|}$.
\end{computation}

\subsubsection{The Full Boundary Divisor $D_{\{1,2,3,4\}}$}

\begin{computation}[$D_{\{1,2,3,4\}}$]
All four points collide. As in the $k=3$ case,
\[
D_{\{1,2,3,4\}}
\cong
\FM_1^{\mathrm{tr}}(\C)\times\FM_4^{\mathrm{red}}(\C),
\]
and this boundary stratum produces the terms involving $m_1$ and $m_4$:
\begin{align*}
&m_1(m_4(a_1,a_2,a_3,a_4)) \\
&+ m_4(m_1(a_1),a_2,a_3,a_4) + (-1)^{|a_1|}\, m_4(a_1,m_1(a_2),a_3,a_4) \\
&+ (-1)^{|a_1|+|a_2|}\, m_4(a_1,a_2,m_1(a_3),a_4) + (-1)^{|a_1|+|a_2|+|a_3|}\, m_4(a_1,a_2,a_3,m_1(a_4)).
\end{align*}
These are the ``$m_1 \leftrightarrow m_4$'' terms in the $n=4$ identity, arising from the Leibniz rule for the differential $Q = m_1$ acting on $m_4$.
\end{computation}

\subsubsection{Assembly: The $n=4$ $A_\infty$ Identity}

\begin{proposition}[$k=4$ Stasheff identity]
\label{prop:k4_identity}
Summing all boundary contributions via Stokes' theorem:
\begin{align*}
0 &= m_1(m_4(a_1,a_2,a_3,a_4)) + \sum_{s=0}^{3} (-1)^{|a_1|+\cdots+|a_s|}\, m_4(a_1,\ldots,m_1(a_{s+1}),\ldots,a_4) \\
&\quad + m_2(m_3(a_1,a_2,a_3),a_4) + (-1)^{|a_1|}\, m_2(a_1,m_3(a_2,a_3,a_4)) \\
&\quad + m_3(m_2(a_1,a_2),a_3,a_4) + (-1)^{|a_1|}\, m_3(a_1,m_2(a_2,a_3),a_4) \\
&\quad + (-1)^{|a_1|+|a_2|}\, m_3(a_1,a_2,m_2(a_3,a_4)),
\end{align*}
with spectral parameters determined by the substitution rules computed above. This is exactly the $n=4$ Stasheff identity~\eqref{eq:ainfty-relation-raw}.

The ten displayed full-chain terms are supported on the six
non-vanishing consecutive boundary divisors:
\begin{itemize}
\item First line (5 terms): $D_{\{1,2,3,4\}}$ contributions ($m_1 \leftrightarrow m_4$);
\item Second line (2 terms): $D_{\{1,2,3\}}$ and $D_{\{2,3,4\}}$ contributions ($m_2 \leftrightarrow m_3$);
\item Third line (3 terms): $D_{\{1,2\}}$, $D_{\{2,3\}}$, $D_{\{3,4\}}$ contributions ($m_3 \leftrightarrow m_2$).
\end{itemize}
Non-consecutive divisors
(\(D_{\{1,3\}}\), \(D_{\{1,4\}}\), \(D_{\{2,4\}}\),
\(D_{\{1,3,4\}}\), \(D_{\{1,2,4\}}\)) contribute zero in the planar
open projection. They are not absent from the closed chiral FM
compactification. If \(m_1=0\), the first line vanishes and the
minimal arity-four identity is the five-term relation
\[
m_2\circ m_3+m_3\circ m_2=0.
\]
The nine-term count belongs instead to the first minimal-model
identity involving \(m_4\), namely the arity-five relation
\[
m_2\circ m_4+m_3\circ m_3+m_4\circ m_2=0.
\]
Corner contributions (codimension-2 strata where two disjoint clusters
collide simultaneously) cancel by the Arnold relations proved in the
next subsection. \qed
\end{proposition}

\begin{remark}[Arity four does not determine \(m_4\) on cohomology]
\label{rem:arity-four-not-m4}
The displayed \(n=4\) identity contains \(m_4\) only through
\(m_1m_4\) and through the terms in which \(m_1\) acts on an input of
\(m_4\). After passing to a minimal model with \(m_1=0\), the
arity-four relation is
\[
m_2\circ m_3+m_3\circ m_2=0.
\]
It is the Jacobi-type coherence for \(m_2\) and \(m_3\); it does not
determine \(m_4\). The first minimal-model Stasheff identity in which
\(m_4\) appears is arity five:
\[
m_2\circ m_4+m_3\circ m_3+m_4\circ m_2=0,
\]
with the usual insertion signs and spectral substitutions. A
chain-level \(m_4\) may be chosen as a homotopy filling the arity-four
source \(m_2\circ m_3+m_3\circ m_2\), but this is an SDR/gauge choice,
not a cohomological determination by the arity-four boundary alone.
\end{remark}

\begin{remark}[Codimension-2 corners in $\FM_4(\C)$]
\label{rem:k4_corners}
In $\FM_4(\C)$, the codimension-2 strata that could potentially contribute are intersections of disjoint boundary divisors, such as $D_{\{1,2\}} \cap D_{\{3,4\}}$ (points 1,2 form one cluster while 3,4 form another, simultaneously). These are genuine codimension-2 corners of the manifold with corners $\FM_4(\C)$. The Arnold cancellation (Theorem~\ref{thm:Arnold_full_proof} below) ensures that the two orderings of taking residues (first $D_{\{1,2\}}$ then $D_{\{3,4\}}$, or vice versa) contribute with opposite signs and cancel. Nested corners such as $D_{\{1,2\}} \cap D_{\{1,2,3\}}$ (where $\{1,2\} \subset \{1,2,3\}$) correspond to hierarchical collisions and are already accounted for as boundary strata of $D_{\{1,2,3\}}$; these produce the terms $m_2(m_2(a_1,a_2),a_3)$ inside $m_3(m_2(\cdot,\cdot),\cdot,\cdot)$ and are part of the codimension-1 analysis.
\end{remark}


\subsection{Arnold--Orlik--Solomon Relations}
\label{subsec:arnold_complete}

\begin{theorem}[Arnold relations on $\FM_n(\C)$]
\label{thm:AOS_forms}
On the configuration space $\Conf_n(\C)$, define the logarithmic $1$-forms
\[
\omega_{ij} := d\log(z_i - z_j) = \frac{dz_i - dz_j}{z_i - z_j}, \qquad 1 \leq i < j \leq n.
\]
These forms satisfy the \emph{Arnold relation}: for any triple $i < j < k$,
\begin{equation}
\label{eq:AOS}
\omega_{ij} \wedge \omega_{jk} + \omega_{jk} \wedge \omega_{ki} + \omega_{ki} \wedge \omega_{ij} = 0.
\end{equation}
Equivalently, writing $\omega_{ji} = -\omega_{ij}$:
\[
\sum_{\text{cyclic}} \omega_{ij} \wedge \omega_{jk} = 0.
\]
\end{theorem}

\begin{proof}
Set $f = z_i - z_j$, $g = z_j - z_k$, so $z_i - z_k = f + g$. Then
\[
\omega_{ij} = \frac{df}{f}, \quad \omega_{jk} = \frac{dg}{g}, \quad \omega_{ik} = \frac{d(f+g)}{f+g}.
\]
The partial fraction identity
\[
\frac{1}{f \cdot g} = \frac{1}{(f+g)} \left(\frac{1}{f} + \frac{1}{g}\right)
\]
implies (multiplying both sides by $df \wedge dg$):
\[
\frac{df}{f} \wedge \frac{dg}{g} = \frac{d(f+g)}{f+g} \wedge \frac{dg}{g} + \frac{df}{f} \wedge \frac{d(f+g)}{f+g}.
\]
To verify this, note that $d(f+g) = df + dg$, so
\begin{align*}
\frac{d(f+g)}{f+g} \wedge \frac{dg}{g} + \frac{df}{f} \wedge \frac{d(f+g)}{f+g}
&= \frac{(df+dg) \wedge dg}{(f+g)g} + \frac{df \wedge (df+dg)}{f(f+g)} \\
&= \frac{df \wedge dg}{(f+g)g} + \frac{df \wedge dg}{f(f+g)} \\
&= df \wedge dg \cdot \frac{1}{f+g}\left(\frac{1}{g} + \frac{1}{f}\right) \\
&= \frac{df \wedge dg}{fg} = \frac{df}{f} \wedge \frac{dg}{g}.
\end{align*}
Rearranging:
\[
\omega_{ij} \wedge \omega_{jk} - \omega_{ik} \wedge \omega_{jk} - \omega_{ij} \wedge \omega_{ik} = 0,
\]
which is~\eqref{eq:AOS} after using $\omega_{ki} = -\omega_{ik}$.
\end{proof}

\begin{corollary}[Arnold is the PVA Jacobi residue identity]
\label{cor:arnold-pva-jacobi-residue}
At arity three, the form identity \eqref{eq:AOS} is the Jacobi
identity for the singular part of \(m_2\).  More precisely, for the
doubly singular three-point form
\(\Omega_3^{\mathrm{sing}\text{-}\mathrm{sing}}(a,b,c)\) of
Lemma~\ref{lem:PVA3_proof}, the three pairwise residues are supported
on
\[
D_{\{1,2\}},\qquad D_{\{2,3\}},\qquad D_{\{1,3\}},
\]
with incidence signs \(+,+,-\).  After the divided-power Borel
transform, and after writing the Jacobi identity in the normal form
\[
J_{\lambda,\mu}(a,b,c)
=
\{a_\lambda\{b_\mu c\}\}
-
(-1)^{(\degree{a}+1)(\degree{b}+1)}
\{b_\mu\{a_\lambda c\}\}
-
\{\{a_\lambda b\}_{\lambda+\mu}c\},
\]
the oriented boundary contributions are
\[
\{a_\lambda\{b_\mu c\}\},\qquad
-\,(-1)^{(\degree{a}+1)(\degree{b}+1)}
\{b_\mu\{a_\lambda c\}\},\qquad
-\{\{a_\lambda b\}_{\lambda+\mu}c\}.
\]
Thus Stokes' theorem on \(\FM_3(\C)\), together with
\(\omega_{12}\wedge\omega_{23}
+\omega_{23}\wedge\omega_{31}
+\omega_{31}\wedge\omega_{12}=0\), gives precisely
Theorem~\ref{thm:Jacobi}.
\end{corollary}

\begin{theorem}[Arnold cancellation at codimension-two corners]
\label{thm:Arnold_full_proof}
Let $S, T \subseteq \{1,\ldots,k\}$ be disjoint subsets with $|S|, |T| \geq 2$. The codimension-2 corner $D_S \cap D_T \subset \FM_k(\C)$ does not contribute to the boundary integral: the two residue orderings cancel exactly.

Precisely, for any closed logarithmic form $\omega \in \Omega^\bullet_{\log}(\FM_k(\C))$ built from wedge products of the $\omega_{ij}$:
\begin{equation}
\label{eq:corner_cancellation}
\int_{\Gamma \cap (D_S)|_{D_T}} \Res_{D_S}\bigl(\Res_{D_T}(\omega)\bigr) + \int_{\Gamma \cap (D_T)|_{D_S}} \Res_{D_T}\bigl(\Res_{D_S}(\omega)\bigr) = 0.
\end{equation}
\end{theorem}

\begin{proof}
\textbf{Local coordinates at the corner.}
Near $D_S \cap D_T$, FM coordinates include two independent radial parameters $\varepsilon_S, \varepsilon_T \geq 0$ (measuring cluster sizes for $S$ and $T$ respectively), along with cluster centers $u_S, u_T$, angular coordinates $\theta_S, \theta_T$, and coordinates for the remaining points. The corner $D_S \cap D_T$ is the locus $\{\varepsilon_S = 0,\, \varepsilon_T = 0\}$.

A logarithmic form near this corner decomposes as
\begin{align*}
\omega &= d\log\varepsilon_S \wedge d\log\varepsilon_T \wedge \gamma_{ST} \\
&\quad + d\log\varepsilon_S \wedge \beta_S + d\log\varepsilon_T \wedge \beta_T + \alpha,
\end{align*}
where $\gamma_{ST}, \beta_S, \beta_T, \alpha$ are smooth at $\varepsilon_S = \varepsilon_T = 0$.

\textbf{Two orderings of residues.}
Taking the residue first at $D_T$ (i.e., setting $\varepsilon_T = 0$ and extracting the coefficient of $d\log\varepsilon_T$):
\[
\Res_{D_T}(\omega) = d\log\varepsilon_S \wedge \gamma_{ST}\big|_{\varepsilon_T=0} + \beta_T\big|_{\varepsilon_T=0}.
\]
Then taking the residue at $D_S$ (extracting $d\log\varepsilon_S$):
\[
\Res_{D_S}\bigl(\Res_{D_T}(\omega)\bigr) = \gamma_{ST}\big|_{\varepsilon_S=\varepsilon_T=0}.
\]

Alternatively, taking residues in the opposite order:
\[
\Res_{D_S}(\omega) = d\log\varepsilon_T \wedge (-1)^1 \gamma_{ST}\big|_{\varepsilon_S=0} + \beta_S\big|_{\varepsilon_S=0},
\]
where the sign $(-1)^1$ arises from anticommuting $d\log\varepsilon_S$ past $d\log\varepsilon_T$ in the first term of $\omega$ (since $\omega = d\log\varepsilon_S \wedge d\log\varepsilon_T \wedge \gamma_{ST} + \cdots = -d\log\varepsilon_T \wedge d\log\varepsilon_S \wedge \gamma_{ST} + \cdots$). Thus
\[
\Res_{D_T}\bigl(\Res_{D_S}(\omega)\bigr) = -\gamma_{ST}\big|_{\varepsilon_S=\varepsilon_T=0}.
\]

\textbf{Orientation signs and cancellation.}
The two factorizations of the corner are:
\begin{itemize}
\item $(D_S)|_{D_T}$: first go to $D_T$ (set $\varepsilon_T = 0$), then within $D_T$ approach $D_S$ (set $\varepsilon_S = 0$). The boundary orientation uses outward normal first for each step. The ambient orientation restricted to the corner is
\[
d\varepsilon_T \wedge d\varepsilon_S \wedge \text{or}(D_S \cap D_T).
\]
\item $(D_T)|_{D_S}$: first go to $D_S$, then approach $D_T$. The orientation is
\[
d\varepsilon_S \wedge d\varepsilon_T \wedge \text{or}(D_S \cap D_T).
\]
\end{itemize}
These two orientations differ by the sign of swapping
$d\varepsilon_S \leftrightarrow d\varepsilon_T$, which is $(-1)$.
With a single orientation fixed on the corner, the incidence numbers of
the two ordered boundary strata differ by this sign. The residue sign
above records the anticommutation of the logarithmic normal forms;
the Stokes contribution uses the residue sign and the incidence sign
as one oriented boundary datum, not as two independently normalised
corner integrals. Thus the two ordered approaches to the same
codimension-$2$ face enter with opposite total signs:
\begin{align*}
&\int_{(D_S)|_{D_T}} \Res_{D_S}(\Res_{D_T}(\omega))
 = \int_{D_S \cap D_T} \gamma_{ST}\big|_{\varepsilon_S=\varepsilon_T=0}, \\
&\int_{(D_T)|_{D_S}} \Res_{D_T}(\Res_{D_S}(\omega))
 = -\int_{D_S \cap D_T} \gamma_{ST}\big|_{\varepsilon_S=\varepsilon_T=0}.
\end{align*}
This is the local form of the identity that each codimension-$2$ face
of a manifold with corners appears twice in the iterated boundary, with
opposite incidence signs. The Stokes orientation formula on a manifold
with corners gives:

The boundary operator $\partial$ on the chain $\FM_k(\C)$ satisfies $\partial^2 = 0$. Writing $\partial = \sum_S \epsilon_S\, D_S$ (sum over codimension-1 faces), the relation $\partial^2 = 0$ means
\[
\sum_{S \neq T} \epsilon_S \epsilon_T \cdot (\text{corner } D_S \cap D_T \text{ with incidence sign}) = 0.
\]
For each pair of disjoint $S, T$, the corner $D_S \cap D_T$ appears exactly twice: once as a boundary face of $D_S$ (the face where $T$ also collapses) and once as a boundary face of $D_T$ (the face where $S$ also collapses). The relation $\partial^2 = 0$ forces these two appearances to carry opposite incidence numbers, which is the content of~\eqref{eq:corner_cancellation}.

This is ultimately a topological fact: on any manifold with corners, $\partial^2 = 0$ on the chain level ensures that corner contributions cancel automatically in Stokes' theorem. The Arnold--Orlik--Solomon relation~\eqref{eq:AOS} provides the \emph{form-theoretic} mechanism by which this cancellation occurs for the specific logarithmic forms built from $\omega_{ij}$. Specifically, if $\omega$ is a product of forms $\omega_{ij}$ and $i \in S$, $j \in T$ with $S \cap T = \varnothing$, then $\omega_{ij}$ contributes a singularity at both $D_S$ and $D_T$. The AOS relation~\eqref{eq:AOS} relates this mixed singularity to purely nested singularities, ensuring that the double residue $\Res_{D_S}(\Res_{D_T}(\omega))$ is controlled by the same partial fraction identity that proves~\eqref{eq:AOS}.
\end{proof}

\begin{example}[Arnold cancellation for $k=4$, $S=\{1,2\}$, $T=\{3,4\}$]
\label{ex:arnold_k4}
Consider the corner $D_{\{1,2\}} \cap D_{\{3,4\}}$ in $\FM_4(\C)$. The weight form involves $\omega_{12}$, $\omega_{23}$, $\omega_{34}$ (for consecutive pairs). Near the corner, both $\varepsilon_{12}$ and $\varepsilon_{34}$ approach zero simultaneously. The double residue $\Res_{D_{\{1,2\}}}(\Res_{D_{\{3,4\}}}(\omega_4))$ extracts the coefficient of $d\log\varepsilon_{12} \wedge d\log\varepsilon_{34}$.

By the AOS relation applied to the triple $(z_1, z_2, z_3)$ and separately to $(z_2, z_3, z_4)$, the logarithmic forms satisfy constraints that prevent an independent double-logarithmic singularity from surviving. Concretely, the form $\omega_{12} \wedge \omega_{34}$ does have a double-log singularity at $D_{\{1,2\}} \cap D_{\{3,4\}}$, but it appears with the same coefficient from both orderings (taking $\Res_{D_{\{1,2\}}}$ first or $\Res_{D_{\{3,4\}}}$ first). The cancellation in~\eqref{eq:corner_cancellation} then follows from $\partial^2 = 0$ as proved above.
\end{example}


\subsection{Rigorous LG Truncation: $m_k = 0$ for $k \geq 4$}
\label{subsec:truncation_complete}

The cubic Landau--Ginzburg vertex leaves no form degree for independent
$k$-ary operations once $k\geq 4$.

\begin{theorem}[Cubic Landau--Ginzburg minimal-model truncation]
\label{thm:LG_truncation_full}
For the Landau--Ginzburg model with cubic superpotential
$W=g\Phi^3/3!$ on $\C\times\R$, the tree-level minimal transferred
\(A_\infty\)-structure on the Jacobi ring \(J(W)=\C[\Phi]/(\Phi^2)\)
satisfies \(m_k=0\) for all \(k\geq4\). Its only non-binary higher
operation is \(m_3(\Phi,\Phi,\Phi)=g\).
\end{theorem}

\begin{proof}
Let \(\pi:\C[\Phi]\to J(W)\) be the Jacobi projection and let \(h\)
be a contraction of the ideal \((\Phi^2)\). The tree-level
homological perturbation formula expresses \(m_n^J\) as a sum over
rooted trees. A vertex of valence \(r\) is labelled by \(W^{(r)}\),
an internal edge is labelled by \(h\), and the root is followed by
\(\pi\). For \(W=g\Phi^3/3!\),
\[
W'''=g,\qquad W^{(r)}=0\quad(r\geq4).
\]
Hence the only one-vertex tree beyond multiplication is the cubic
corolla, and it gives \(m_3(\Phi,\Phi,\Phi)=g\). A tree with at least
two cubic vertices has an internal cubic edge; after multiplication
along this edge its root contribution lies in \((\Phi^2)\). The
projection \(\pi\) kills this ideal. Therefore every rooted tree with
\(n\geq4\) vanishes in the minimal transferred Jacobi model.

This is a tree-level statement. Loop graphs live in the renormalised BV
theory before projection to \(J(W)\); they are not constrained by this
minimal-model argument.
\end{proof}

\begin{remark}[Comparison with the quartic case]
For a quartic superpotential $W = \frac{g}{4!}\Phi^4$, each vertex is $4$-valent. A tree with $k$ external legs has $V = (k-2)/2$ internal vertices (requiring $k$ even for non-zero trees) and $E = V - 1 = (k-4)/2$ internal edges. The holomorphic form deficit is $V - E = 1$, and the same vanishing argument applies for $k \geq 5$ (with $m_4 \neq 0$ being the first non-trivial higher operation). More generally, for an $r$-valent vertex, the operation $m_r$ is the first non-trivial higher operation, and $m_k = 0$ for $k > r$ at tree level.
\end{remark}
